\documentclass[../../../include/open-logic-section]{subfiles}

\begin{document}

\olfileid[id]{his}{bio}{rus}

\olsection{Bertrand Russell}

\olphoto{russell-bertrand}{Bertrand Russell}

Bertrand Russell dipuji sebagai salah seorang pendiri filsafat analitik
modern. Lahir pada 18 Mei 1872, Russell tidak hanya dikenal karena karyanya
dalam filsafat dan logika, tetapi juga menulis banyak buku populer dalam
berbagai bidang. Sepanjang hidupnya ia juga merupakan aktivis politik yang
gigih.

Russell lahir di Trellech, Monmouthshire, Wales. Orang tuanya merupakan
anggota bangsawan Britania. Mereka adalah pemikir bebas dan bahkan
bersahabat dengan kaum radikal di Boston pada masa itu. Sayangnya, orang tua
Russell meninggal ketika ia masih kecil, dan Russell kemudian tinggal
bersama kakek-neneknya. Di sana ia dibesarkan secara religius (sesuatu yang
ingin dihindari orang tuanya dengan segala cara). Neneknya sangat ketat
dalam segala perkara moral. Semasa remaja ia terutama bersekolah di rumah
dengan tutor privat.

Pengaruh Russell dalam filsafat analitik, terutama logika, sangat besar. Ia
mempelajari matematika dan filsafat di Trinity College, Cambridge, tempat ia
dipengaruhi oleh matematikawan dan filsuf Alfred North Whitehead. Pada tahun
1910, Russell dan Whitehead menerbitkan jilid pertama \emph{Principia
Mathematica}, tempat mereka memperjuangkan pandangan bahwa matematika dapat
direduksi menjadi logika. Ia kemudian menerbitkan ratusan buku, esai, dan
pamflet politik. Pada tahun 1950, ia meraih Hadiah Nobel Sastra.

Russell terlibat sangat dalam dalam politik dan aktivisme sosial. Selama
Perang Dunia~I ia ditangkap dan dipenjara selama enam bulan karena kegiatan
dan protes pasifis. Di penjara, ia dapat menulis dan membaca, dan mengaku
menganggap pengalaman itu ``cukup menyenangkan.'' Ia tetap menjadi pasifis
sepanjang hidupnya dan kembali dipenjara karena menghadiri unjuk rasa
pelucutan senjata nuklir pada tahun 1961. Ia juga selamat dari kecelakaan
pesawat pada tahun 1948, ketika satu-satunya orang yang selamat adalah mereka
yang duduk di bagian merokok. Karena itu, Russell mengatakan bahwa ia
berutang nyawa kepada rokok. Russell menikah empat kali, tetapi dikenal
kerap menjalin hubungan di luar pernikahan. Ia meninggal pada 2 Februari
1970 dalam usia 97 tahun di Penrhyndeudraeth, Wales.

\begin{reading}
Russell menulis autobiografi dalam tiga bagian yang merentang kehidupannya
dari tahun 1872--1967 \citep{Russell1967,Russell1968,Russell1969}. Bertrand
Russell Research Centre di McMaster University merupakan tempat arsip
Bertrand Russell disimpan. Lihat situs web mereka di \citet{Duncan2015}
untuk informasi tentang jilid-jilid kumpulan karyanya (termasuk indeks yang
dapat ditelusuri) dan proyek-proyek kearsipan. Makalah Russell \emph{On
  Denoting} \citep{Russell1905} merupakan karya klasik filsafat analitik abad
ke-20.

Entri mengenai Russell dalam Stanford Encyclopedia of Philosophy
\citep{Irvine2015} memuat rekaman suara Russell yang berbicara tentang hasrat
dan teori politik. Banyak wawancara video dengan Russell tersedia secara
daring. Sebagai contoh, untuk melihatnya berbicara tentang merokok dan
keterlibatannya dalam kecelakaan pesawat, lihat \citet{RussellND}. Beberapa
karya Russell, termasuk \emph{Introduction to Mathematical Philosophy},
tersedia sebagai buku audio gratis di \citet{LibriVoxND}.
\end{reading}

\end{document}
